\nonstopmode

\newif\ifcomments
\commentstrue

%% For double-blind review submission
%%\documentclass[acmsmall%acmlarge%,review,anonymous
%%]{acmart}\settopmatter{printfolios=true}
%% For single-blind review submission
\documentclass[svgnames,dvipsnames,table,review%,anonymous
]{jfp-epi}
\settopmatter{printfolios=true}

%% For final camera-ready submission
%\documentclass[dvipsnames]{jfp-epi}\settopmatter{}

\jfpVolume{99}
\jfpArticle{13}
\jfpDOI{10.46298/jfp.01234}
\jfpYear{2027}

\received[Submitted]{October 2026}
%\received[accepted]{April 2089}

\usepackage[utf8]{inputenc}
\usepackage[verbose]{newunicodechar}
\usepackage[capitalise]{cleveref} % capitalise to please Andreas :)
% \usepackage{thmtools} % breaks many links. But not needed
\Crefname{example}{Example}{Examples}
\Crefname{definition}{Definition}{Definitions}
\Crefname{law}{Law}{Laws}
\Crefname{lemma}{Lemma}{Lemmas}
\Crefname{corollary}{Corollary}{Corollaries}
\usepackage[textsize=tiny,disable]{todonotes}
\usepackage[misc]{ifsym}

% Remove todo-notes in final version
% \renewcommand{\todo}[1]{}

\ifcomments
\newcommand\andreas[1]{\todo[backgroundcolor=green!10]{AA: #1}}
\newcommand\koen[1]{\todo[backgroundcolor=blue!10]{KC: #1}}
\newcommand\christian[1]{\todo[backgroundcolor=green!10]{CS: #1}}
\else
\newcommand\andreas[1]{}
\newcommand\koen[1]{}
\newcommand\christian[1]{}
\fi

%\usepackage{tikz}
%\usetikzlibrary{matrix}

% \newunicodechar{↝}{\ensuremath{\leadsto}}
%\newunicodechar{⁺}{^+}
%\newunicodechar{₁}{_1}
%\newunicodechar{₂}{_2}
\DeclareMathAlphabet{\mymathds}{U}{BOONDOX-ds}{m}{n}
\newcommand\llbrace{\mathopen{{\{\kern-0.24em|}}}    % opening {{ value delimiter
\newcommand\rrbrace{\mathclose{{|\kern-0.24em\}}}}   % closing }} value delimiter
\newcommand{\bidirmultimap}{\mathrel{\ooalign{$\mkern6mu\multimap$\cr$\reflectbox{\(\multimap\)}$\cr}}}
\newunicodechar{ }{{\,}}
\newunicodechar{¬}{\ensuremath{\neg}} %
\newunicodechar{²}{^2}
\newunicodechar{³}{^3}
\newunicodechar{·}{\ensuremath{\cdot}}
\newunicodechar{¹}{^1}
\newunicodechar{×}{\ensuremath{\times}} %
\newunicodechar{÷}{\ensuremath{\div}} %
\newunicodechar{ʳ}{^r}
\newunicodechar{ˡ}{^l}
\newunicodechar{Γ}{\Gamma}   %
\newunicodechar{Δ}{\ensuremath{\Delta}} %
\newunicodechar{Η}{\ensuremath{\textrm{H}}} %
\newunicodechar{Θ}{\ensuremath{\Theta}} %
\newunicodechar{Λ}{\ensuremath{\Lambda}} %
\newunicodechar{Ξ}{\ensuremath{\Xi}} %
\newunicodechar{Π}{\ensuremath{\Pi}}   %
\newunicodechar{Σ}{\ensuremath{\Sigma}} %
\newunicodechar{Φ}{\ensuremath{\Phi}} %
\newunicodechar{Ψ}{\ensuremath{\Psi}} %
\newunicodechar{Ω}{\ensuremath{\Omega}} %
\newunicodechar{α}{\ensuremath{\mathnormal{\alpha}}} %
\newunicodechar{β}{\ensuremath{\beta}} %
\newunicodechar{γ}{\ensuremath{\mathnormal{\gamma}}} %
\newunicodechar{δ}{\ensuremath{\mathnormal{\delta}}} %
\newunicodechar{ε}{\ensuremath{\mathnormal{\varepsilon}}} %
\newunicodechar{ζ}{\ensuremath{\mathnormal{\zeta}}} %
\newunicodechar{η}{\ensuremath{\mathnormal{\eta}}} %
\newunicodechar{θ}{\ensuremath{\mathnormal{\theta}}} %
\newunicodechar{ι}{\ensuremath{\mathnormal{\iota}}} %
\newunicodechar{κ}{\ensuremath{\mathnormal{\kappa}}} %
\newunicodechar{λ}{\ensuremath{\mathnormal{\lambda}}} %
\newunicodechar{μ}{\ensuremath{\mathnormal{\mu}}} %
\newunicodechar{ν}{\ensuremath{\mathnormal{\mu}}} %
\newunicodechar{ξ}{\ensuremath{\mathnormal{\xi}}} %
\newunicodechar{π}{\ensuremath{\mathnormal{\pi}}} %
\newunicodechar{ρ}{\ensuremath{\mathnormal{\rho}}} %
\newunicodechar{σ}{\ensuremath{\mathnormal{\sigma}}} %
\newunicodechar{τ}{\ensuremath{\mathnormal{\tau}}} %
\newunicodechar{φ}{\ensuremath{\mathnormal{\varphi}}} %
\newunicodechar{χ}{\ensuremath{\mathnormal{\chi}}} %
\newunicodechar{ψ}{\ensuremath{\mathnormal{\psi}}} %
\newunicodechar{ω}{\ensuremath{\mathnormal{\omega}}} %
\newunicodechar{ϕ}{\ensuremath{\mathnormal{\phi}}} %
\newunicodechar{ϵ}{\ensuremath{\mathnormal{\epsilon}}} %
\newunicodechar{ᵏ}{\ensuremath{^k}}
\newunicodechar{ᵢ}{\ensuremath{_i}} %
\newunicodechar{•}{\ensuremath{\cdot}} %
\newunicodechar{ }{{\,}}
\newunicodechar{′}{\ensuremath{^\prime}}  %
\newunicodechar{″}{\ensuremath{^\second}} %
\newunicodechar{‴}{\ensuremath{^\third}}  %
\newunicodechar{ⁱ}{\ensuremath{^i}}
\newunicodechar{⁵}{\ensuremath{^5}}
\newunicodechar{⁺}{\ensuremath{^+}} %%
\newunicodechar{⁻}{\ensuremath{^-}} %%
\newunicodechar{ⁿ}{\ensuremath{^n}}
\newunicodechar{₀}{\ensuremath{_0}} %
\newunicodechar{₁}{\ensuremath{_1}}
\newunicodechar{₂}{\ensuremath{_2}}
\newunicodechar{₃}{\ensuremath{_3}}
\newunicodechar{₊}{\ensuremath{_+}} %%
\newunicodechar{₋}{\ensuremath{_-}} %%
\newunicodechar{ₙ}{\ensuremath{_n}}
\newunicodechar{ₚ}{\ensuremath{_p}} %
\newunicodechar{ℂ}{\ensuremath{\mathds{C}}} %
\newunicodechar{ℕ}{\mathbb{N}} %
\newunicodechar{ℝ}{\ensuremath{\mathbb{R}}} %
\newunicodechar{ℤ}{\ensuremath{\mathbb{Z}}} %
\newunicodechar{ℳ}{\mathcal{M}}
\newunicodechar{⅋}{\ensuremath{\parr}} %
\newunicodechar{←}{\ensuremath{\leftarrow}} %
\newunicodechar{↑}{\ensuremath{\uparrow}} %
\newunicodechar{↓}{\ensuremath{\downarrow}} %
\newunicodechar{→}{\ensuremath{\rightarrow}} %
\newunicodechar{↔}{\ensuremath{\leftrightarrow}} %
\newunicodechar{↖}{\nwarrow} %
\newunicodechar{↗}{\nearrow} %
\newunicodechar{↝}{\ensuremath{\rightsquigarrow}}
\newunicodechar{↦}{\ensuremath{\mapsto}}
\newunicodechar{↭}{\ensuremath{\leftrightsquigarrow}} %
\newunicodechar{↾}{\ensuremath{\upharpoonright}} %
\newunicodechar{↿}{\ensuremath{\upharpoonleft}} %
\newunicodechar{⇆}{\ensuremath{\leftrightarrows}} %
\newunicodechar{⇐}{\ensuremath{\Leftarrow}} %
\newunicodechar{⇒}{\ensuremath{\Rightarrow}} %
\newunicodechar{⇔}{\ensuremath{\Leftrightarrow}} %
\newunicodechar{∀}{\ensuremath{\forall}}   %
\newunicodechar{∃}{\ensuremath{\exists}} %
\newunicodechar{∅}{\ensuremath{\varnothing}} %
\newunicodechar{∈}{\ensuremath{\in}}
\newunicodechar{∉}{\ensuremath{\not\in}} %
\newunicodechar{∋}{\ensuremath{\ni}}  %
\newunicodechar{∎}{\ensuremath{\qed}}
\newunicodechar{∏}{\prod}
\newunicodechar{∑}{\sum}
\newunicodechar{∗}{\ensuremath{\ast}} %
\newunicodechar{∘}{\ensuremath{\circ}} %
\newunicodechar{∙}{\ensuremath{\bullet}} %
\newunicodechar{∙}{\ensuremath{\cdot}}
\newunicodechar{∞}{\ensuremath{\infty}} %
\newunicodechar{∣}{\ensuremath{\mid}} %
\newunicodechar{∧}{\ensuremath{\wedge}}%
\newunicodechar{∨}{\ensuremath{\vee}}%
\newunicodechar{∩}{\ensuremath{\cap}} %
\newunicodechar{∪}{\ensuremath{\cup}} %
\newunicodechar{∫}{\int}
\newunicodechar{∷}{::} %
\newunicodechar{∼}{\ensuremath{\sim}} %
\newunicodechar{≃}{\ensuremath{\simeq}} %
\newunicodechar{≅}{\ensuremath{\cong}} %
\newunicodechar{≈}{\ensuremath{\approx}} %
\newunicodechar{≔}{:=} %
\newunicodechar{≜}{\ensuremath{\stackrel{\scriptscriptstyle {\triangle}}{=}}}
\newunicodechar{≜}{\triangleq}
\newunicodechar{≝}{\ensuremath{\stackrel{\scriptscriptstyle {\text{def}}}{=}}}
\newunicodechar{≟}{\ensuremath{\stackrel {_\text{\textbf{?}}}{\text{\textbf{=}}\negthickspace\negthickspace\text{\textbf{=}}}}}
\newunicodechar{≠}{\ensuremath{\neq}}%
\newunicodechar{≡}{\ensuremath{\equiv}}%
\newunicodechar{≤}{\ensuremath{\leq}} %
\newunicodechar{≥}{\ensuremath{\geq}} %
\newunicodechar{≺}{\ensuremath{\prec}}
\newunicodechar{≼}{\ensuremath{\preceq}}
\newunicodechar{⊂}{\ensuremath{\subset}} %
\newunicodechar{⊃}{\ensuremath{\supset}} %
\newunicodechar{⊆}{\ensuremath{\subseteq}} %
\newunicodechar{⊇}{\ensuremath{\supseteq}} %
\newunicodechar{⊎}{\ensuremath{\uplus}} %
\newunicodechar{⊑}{\ensuremath{\sqsubseteq}} %
\newunicodechar{⊒}{\ensuremath{\sqsupseteq}} %
\newunicodechar{⊓}{\ensuremath{\sqcap}} %
\newunicodechar{⊔}{\ensuremath{\sqcup}} %
\newunicodechar{⊕}{\ensuremath{\oplus}} %
\newunicodechar{⊗}{\ensuremath{\otimes}} %
\newunicodechar{⊞}{\ensuremath{\boxplus}} %
\newunicodechar{⊟}{\ensuremath{\boxminus}} %
\newunicodechar{⊡}{\ensuremath{\boxdot}} %
\newunicodechar{⊢}{\ensuremath{\vdash}} %
\newunicodechar{⊣}{\ensuremath{\dashv}}
\newunicodechar{⊤}{\ensuremath{\top}}%
\newunicodechar{⊥}{\ensuremath{\bot}} % bottom
\newunicodechar{⊧}{\vDash} %
\newunicodechar{⊨}{\vDash} %
\newunicodechar{⊩}{\Vdash}
\newunicodechar{⊸}{\ensuremath{\multimap}} %
\newunicodechar{⋁}{\ensuremath{\bigvee}}
\newunicodechar{⋃}{\ensuremath{\bigcup}} %
\newunicodechar{⋄}{\ensuremath{\diamond}} %
\newunicodechar{⋆}{\ensuremath{\star}} %
\newunicodechar{⋮}{\ensuremath{\vdots}} %
\newunicodechar{⋯}{\ensuremath{\cdots}} %
\newunicodechar{─}{---}
\newunicodechar{□}{\ensuremath{\Box}} %
\newunicodechar{▹}{\ensuremath{\rhd}} %
\newunicodechar{◂}{\ensuremath{\triangleleft}}
\newunicodechar{▴}{\ensuremath{\mathbin{\triangle}}}
\newunicodechar{◃}{\triangleleft{}} %
\newunicodechar{◆}{\blacklozenge} %
\newunicodechar{◇}{\ensuremath{\Diamond}} %
\newunicodechar{◽}{\ensuremath{\square}} %
\newunicodechar{★}{\ensuremath{\star}}   %
\newunicodechar{♢}{\ensuremath{\diamondsuit}} %
\newunicodechar{♭}{\ensuremath{\flat}} %
\newunicodechar{♯}{\ensuremath{\sharp}} %
\newunicodechar{✓}{\ensuremath{\checkmark}} %
\newunicodechar{⟂}{\ensuremath{\bot}} % PERPENDICULAR
\newunicodechar{⟦}{\ensuremath{\llbracket}}
\newunicodechar{⟧}{\ensuremath{\rrbracket}}
\newunicodechar{⟨}{\ensuremath{\langle}} %
\newunicodechar{⟩}{\ensuremath{\rangle}} %
\newunicodechar{⟶}{\ensuremath{\longrightarrow}}
\newunicodechar{⟷}{\longleftrightarrow}
\newunicodechar{⟹}{\ensuremath{\Longrightarrow}} %
\newunicodechar{⦃}{\ensuremath{\llbrace}} %
\newunicodechar{⦄}{\ensuremath{\rrbrace}} %
\newunicodechar{⧟}{\ensuremath{\bidirmultimap}} %
\newunicodechar{⨁}{\bigoplus} %
\newunicodechar{⨂}{\bigotimes} %
\newunicodechar{⫪}{\ensuremath{\top\kern-1.3ex\top}}
\newunicodechar{ⱼ}{_j} %
\newunicodechar{𝒟}{\ensuremath{\mathcal{D}}} %
\newunicodechar{𝒢}{\ensuremath{\mathcal{G}}}
\newunicodechar{𝒦}{\ensuremath{\mathcal{K}}} %
\newunicodechar{𝔸}{\ensuremath{\mathds{A}}} %
\newunicodechar{𝔹}{\ensuremath{\mathds{B}}} %
\newunicodechar{𝟘}{\ensuremath{\mymathds{0}}} %
\newunicodechar{𝟙}{\ensuremath{\mymathds{1}}} %


% For having two versions of a paper in one file.
% Stuff that does not fit into the short version can be encosed in \LONGVERSION{...}

\ifdefined\LONGVERSION
\newcommand{\SHORTVERSION}[1]{}
\else
% short version:
\newcommand{\LONGVERSION}[1]{}
\newcommand{\SHORTVERSION}[1]{#1}
% % long version:
% \newcommand{\LONGVERSION}[1]{#1}
% \newcommand{\SHORTVERSION}[1]{}
% \newcommand{\SHORTVERSION}[1]{BEGIN~SHORT\ #1 \ END~SHORT}
\fi
\newcommand{\LONGSHORT}[2]{\LONGVERSION{#1}\SHORTVERSION{#2}}
\newcommand{\SHORTLONG}[2]{\SHORTVERSION{#1}\LONGVERSION{#2}}
\newcommand{\EXTENDED}[1]{}
\newwrite\agdarefs
\immediate\openout\agdarefs=References.agda

% \newcommand\agdaref[2]{\href{file:/home/jyp/repo/modal/agda/html/#1.html\##2}{✓}}
\def\myhash{\#}
\write\agdarefs{{-\string# OPTIONS -W all --warning=error \string#-}}
\write\agdarefs{module References where}
\write\agdarefs{}
\newcommand\agdacomment[1]{\write\agdarefs{-- #1}}
\newcommand\agdablank{\write\agdarefs{}}
\NewDocumentCommand\agdaref{o m o m}{%
  \href{html/#2.html\#\IfValueT{#3}{#3.}#4}{✓}%
  \IfValueT{#1}{\agdacomment{#1}}%
  \write\agdarefs{import #2}%
  \write\agdarefs{open #2\IfValueT{#3}{.#3} using (#4)} %
  \write\agdarefs{}%
}
\newcommand\magdaref[2]{\text{\agdaref{#1}{#2}}}

\newcommand{\NOVERSION}[1]{}  %% currently not included
\newcommand\seeExtendedMaterial{(See extended material)}

\SHORTVERSION{
\newcommand\proofInSupp{The proof can be found in the long version of
  the paper, provided as supplementary material.
  \renewcommand\proofInSupp{\hfill\seeExtendedMaterial} }
}
\LONGVERSION{
\newcommand\proofInSupp{}
}

\newenvironment{definition*}{}{} % if this should not be emphasized
\makeatletter
\newenvironment{proof*}[1][\proofname]{\par
  \normalfont \topsep6\p@\@plus6\p@\relax
  \trivlist
  \item[\@proofindent\hskip\labelsep
        {\@proofnamefont #1\@addpunct{.}}]\ignorespaces
}{%
}
\makeatother
%\usepackage[right]{showlabels}\renewcommand{\showlabelfont}{\small\ttfamily\color{gray}}

%\usepackage{prooftree}
%\usepackage{dsfont}
%\usepackage{mathpartir}
%\usepackage{mathtools} %mathrlap, prescript
% \usepackage[all]{xy}
%\usepackage{colortbl}

% \usepackage[cal=boondoxo]{mathalfa}
%\usepackage{calc} % \widthof

\usepackage{latex/agda}

% NO EFFECT:
% \renewcommand{\AgdaOperator}    [1]
%     {\AgdaNoSpaceMath{\AgdaFontStyle{\textcolor{AgdaOperator}{\mathbf{#1}}}}}

%% Note: Authors migrating a paper from PACMPL format to traditional
%% SIGPLAN proceedings format should change 'acmlarge' to
%% 'sigplan,10pt'.



%% Some recommended packages.
% \usepackage{booktabs}   %% For formal tables:
%                         %% http://ctan.org/pkg/booktabs
% \usepackage{subcaption} %% For complex figures with subfigures/subcaptions
%                         %% http://ctan.org/pkg/subcaption

\newtheorem{remark}{Remark}

%% Bibliography style
\bibliographystyle{ACM-Reference-Format}
%% Citation style
% \citestyle{acmauthoryear}   %% For author/year citations
% \citestyle{numbers,sort&compress}

\input{textMacros}

% Macros for this paper
\newcommand{\rot}{\ensuremath{\mathord{\mathsf{r}}}}
\newcommand{\tail}{\ensuremath{\mathord{\mathsf{t}}}}
\newcommand{\tC}{\ensuremath{\mathsf{C}}}
\newcommand{\tR}{\ensuremath{\mathsf{R}}}
\newcommand{\tT}{\ensuremath{\mathsf{T}}}
\renewcommand{\Cat}[2]{\ensuremath{\tC_{#1,#2}}}
\newcommand{\Rot}[1]{\ensuremath{\tR_{#1}}}
\renewcommand{\Tail}[1]{\ensuremath{\tT_{#1}}}
\newcommand{\vms}{\ensuremath{\mathit{ms}}}
\newcommand{\vts}{\ensuremath{\mathit{ts}}}
\newcommand{\vtsp}{\ensuremath{\mathit{ts'}}}



\begin{document}

%% Title information
\title{%
  Binary trees are constant-time catenable queues
  \LONGVERSION{ (Extended~Version)}}
%\titlenote{with title note}             %% \titlenote is optional;
                                        %% can be repeated if necessary;
                                        %% contents suppressed with 'anonymous'
%\subtitle{Subtitle}                     %% \subtitle is optional
%\subtitlenote{with subtitle note}       %% \subtitlenote is optional;
                                        %% can be repeated if necessary;
                                        %% contents suppressed with 'anonymous'


%% Author information
%% Contents and number of authors suppressed with 'anonymous'.
%% Each author should be introduced by \author, followed by
%% \authornote (optional), \orcid (optional), \affiliation, and
%% \email.
%% An author may have multiple affiliations and/or emails; repeat the
%% appropriate command.
%% Many elements are not rendered, but should be provided for metadata
%% extraction tools.

\author{Andreas Abel}
%\authornote{with author1 note}          %% \authornote is optional;
                                        %% can be repeated if necessary
\orcid{0000-0003-0420-4492}             %% \orcid is optional
\affiliation{
  \department{Department of Computer Science and Engineering}
  \institution{University of Gothenburg and Chalmers University of Technology}
  \city{Göteborg}
%  \state{State1}
  \country{Sweden}
}
\email{andreas.abel@gu.se}          %% \email is recommended


\author{Koen Claessen}
%\authornote{with author1 note}          %% \authornote is optional;
                                        %% can be repeated if necessary
\orcid{}             %% \orcid is optional
\affiliation{
  \department{Department of Computer Science and Engineering}
  \institution{University of Gothenburg and Chalmers University of Technology}
  \city{Göteborg}
%  \state{State1}
  \country{Sweden}
}
\email{}          %% \email is recommended

\author{Christian Sattler}
%\authornote{with author1 note}          %% \authornote is optional;
                                        %% can be repeated if necessary
\orcid{}             %% \orcid is optional
\affiliation{
  \department{Department of Computer Science and Engineering}
  \institution{University of Gothenburg and Chalmers University of Technology}
  \city{Göteborg}
%  \state{State1}
  \country{Sweden}
}
\email{}          %% \email is recommended



%% Paper note
%% The \thanks command may be used to create a "paper note" ---
%% similar to a title note or an author note, but not explicitly
%% associated with a particular element.  It will appear immediately
%% above the permission/copyright statement.
%\thanks{with paper note}                %% \thanks is optional
                                        %% can be repeated if necesary
                                        %% contents suppressed with 'anonymous'


%% Abstract
%% Note: \begin{abstract}...\end{abstract} environment must come
%% before \maketitle command
\begin{abstract}
  Binary trees are constant-time catenable queues.
  The paper is backed by a mechanisation
  in Agda.
\end{abstract}


%% \maketitle
%% Note: \maketitle command must come after title commands, author
%% commands, abstract environment, Computing Classification System
%% environment and commands, and keywords command.
\maketitle


% \fleq % nyarlathotep
%\fleq  %% Flushleft equations
%\cneq  %% Center equations


\section{Introduction}\label{sec:intro}

Binary trees are constant-time catenable queues.

% https://hackage.haskell.org/package/alist

\section{Explicit rotation}\label{sec:rotation}

%\input{latex/Code}

We consider a modified API that makes tree rotation a primitive operation, since this simplifies the amortization proofs.
Since the contents of the queue does not matter for complexity considerations, only its \emph{structure}, we drop the labels from the tree definition and only consider the cases of an (empty) leaf \(ε\) and a binary node constructor \(\_∙\_\).
Tree rotation \(\_\rot\) is a partial function defined by \(((t₁ ∙ t₂) ∙ t₃)\rot = t₁ ∙ (t₂ ∙ t₃)\).
Dropping from the front of the queue is represented by the partial function ``tail'' \(\_\tail\) defined by \( (ε ∙ t)\tail = t \).
The idea is that the tail of a queue can only be taken if its tree representation has been maximally rotated at the root and thus is of the form \( ε ∙ t \).

\subsection{The tree game}

To study the amortized complexity, we need to talk about arbitrary sequences of tree operations.
We conceptualize such sequences as moves in a game where the arena is a multiset of trees which we call a \emph{forest}.
Each move takes out some trees from the forest, applies the corresponding operation on them and puts the result back into the forest.
For simplicity we represent the forest as a non-empty list so that we can identify particular trees by their index in the list.  Note that taking out trees from the forest and putting a tree back (which we do at the front of the list) typically changes indices of tree in the forest.
The moves $m$ we are considering are:
\[
\begin{array}{l@{\qquad}l}
  \Cat i j & \mbox{concatenate trees $i$ and $j$} \\
  \Rot i   & \mbox{rotate tree $i$, if possible} \\
  \Tail i  & \mbox{take tail of tree $i$, if possible} \\
\end{array}
\]
We write $\vts \, m = \vtsp$ if move $m$ can be applied to forest $\vts$, yielding forest $\vts'$.
For instance, $(t₀, t₁, t₂, t₃) \Cat 1 2 = (t₁ ∙ t₃, t₀, t₂)$ since after taking out $t₁$ the second index, $2$, refers to $t₃$ in the remaining forest $(t₀, t₂, t₃)$.
We may omit indices if they are $0$.
E.g. $(ε, ε, ε)\tC\tC = ((ε ∙ ε) ∙ ε)$ and $(ε, ε, ε)\tC\tC\tR\tT\tT = (ε)$.

It is clear that the only terminal forest is $(ε)$; no move is applicable here, but any other forest has at least one move.  Any tree that is not $ε$ can either be rotated or truncated to its tail, and if the forest has more than one tree, concatenation is possible.
A forest of $n$ leaf-trees $\varepsilon$ is called an \emph{initial} forest and written $\varepsilon^n$.
We are most interested in move sequences starting in initial forests.

To establish amortized complexity we are interested in the longest move sequences starting from a given forest, and the amount of $\tC$, $\tT$ and $\tR$ moves in those sequences.
It is obvious that the number of $\tC$ moves in a longest sequence is the number of trees in the forest minus one, since each such move reduces the number of trees by one.
Since each $\tT$ move removes one leaf but the last one, the number of $\tT$ moves is at most the number of leaves in the forest minus one.
The maximal number of $\tR$ moves is the quantity we are most interested in.
If it is linearly bounded by the size of the forest, we can execute rotations for free, yielding amortized constant time for queue extraction which in the general case consists of a number of rotations and then a tail operation.
We will prove that in move sequences starting in an inital forest $\varepsilon^{n+1}$ the number of rotations will be at most $4n$.

\subsection{Resourced execution}

To amortize the rotations, we raise some fees $k_\tC$ and $k_\tT$ that have to be paid for concatenation and tail operations.
Those are stored on a virtual bank account $k$ (attached to a forest as subscript $\vts_k$) so that we can pay a unit cost of $1$ for each rotation.
For uniformity, let us set $k_\tR = -1$ and refine move execution by the bank account:
\[
  \vts_k \, m = \vtsp_{k'} \mbox{ if } k' = k + k_m \geq 0 \mbox{ and } \vts \, m = \vtsp
  .
\]
The goal is to show that there are constants $k_\tC$ and $k_\tT$ such that for each legal move sequence $\vts \, \vms = \vtsp$ there is some initial resource $k$ and some leftover resource $k' \geq 0$ such that $\vts_k \, \vms = \vtsp_{k'}$.
Moreover, for an initial forest $\vts = \varepsilon^n$ the resourced execution should succeed without initial resources, so $k=0$.

One solution is $k_\tC = k_\tT = 2$, and we prove its correctness via the potential method.
To this end, we define the following potential functions on forests and trees:
\[
\begin{array}{lll}
\Phi(t_0,\dots,t_n) & = & \sum_{i=0..n} \Phi(t_i) \\
\Phi(\varepsilon)   & = & \Phi_l(\varepsilon) = \Phi_r(\varepsilon) = 0 \\
\Phi(t₁ ∙ t₂)       & = & \max (\Phi_l(t₁)+1, \Phi_r(t₂)-1) \\
\Phi_l(t₁ ∙ t₂)     & = & \max (\Phi_l(t₁)+1, \Phi_r(t₂)) \\
\Phi_r(t₁ ∙ t₂)     & = & \max (\Phi_r(t₁)+1, \Phi_r(t₂)-1) \\
\end{array}
\]
Clearly, all of these functions yield non-negative integers.
The various potential functions on trees are related as follows:
\begin{lemma} \hfill
  \label{lem:phi}
  \begin{enumerate}
  \item $\Phi_r(t) \leq \Phi_l(t) \leq 1 + \Phi_r(t)$.
  \item $\Phi_r(t) \leq \Phi(t) \leq 1 + \Phi_r(t)$.
  \item $\Phi_l(t) \leq 1 + \Phi(t)$.
  \end{enumerate}
\end{lemma}
\begin{proof}
  Easy arithmetical proofs by induction on $t$.
\end{proof}

\begin{theorem}
  If $\vts \, \vms = \vtsp$ and $k \geq \Phi(\vts)$ then $\vts_k \, \vms = \vtsp_{k'}$ for some $k' \geq \Phi(\vtsp)$.
\end{theorem}
\begin{proof*}
  By induction on the moves.  It is sufficient to consider single moves.
  \begin{caselist}

  \nextcase $m = \Cat i j$.
  Then there is some partition of $\vts$ into two trees $t_1$ and $t_2$ and some rest $\vtsp$.
  We have to show that for $k \geq \Phi(t_1) + \Phi(t_2) + \Phi(\vtsp)$
  there is some $k' \geq \Phi(t_1 ∙ t_2) + \Phi(\vtsp)$ such that $k + 2 = k'$.
  It is sufficient to show that $\Phi(t_1) + \Phi(t_2) + 2 \geq \Phi(t_1 ∙ t_2) = \max (\Phi_l(t₁)+1, \Phi_r(t₂)-1)$.
  This follows from $\Phi(t_1) + 2 \geq \Phi_l(t₁) + 1$
  and $\Phi(t_2) + 2 \geq \Phi_r(t₂)-1$, both implied by \Cref{lem:phi}.

  \nextcase $m = \Tail i$.
  Then forest $\vts$ is partitioned into a tree $t_i = ε ∙ t$ and a remainder $\vtsp$.
  If is sufficient to show that $\Phi(ε ∙ t) + 2 = \max(1, \Phi_r(t)) + 2 \geq \Phi(t)$.
  This follows from  \Cref{lem:phi}, $\Phi_r(t)) + 1 \geq \Phi(t)$.

  \nextcase $m = \Rot i$.
  Then forest $\vts$ is partitioned into a tree $(t₁ ∙ t₂) ∙ t₃$ and a remainder $\vtsp$.
  It is sufficient to show that
  $\Phi((t₁ ∙ t₂) ∙ t₃) \geq \Phi(t₁ ∙ (t₂ ∙ t₃)) + 1$
  but both sides are equal to $\max(\Phi_l(t₁)+2, \Phi_r(t₂)+1, \Phi_r(t₃)-1)$.
  \qed

  \end{caselist}
\end{proof*}
\begin{corollary}
  If $\varepsilon^n \, \vms = \vtsp$ then $\varepsilon^n_0 \, \vms = \vtsp_{k'}$ for some $k' \geq 0$.
\end{corollary}


\subsection{Optimality}

We have shown that $k_\tC = k_\tT = 2$ is sufficient to amortize all rotations, but is this budget optimal?
In the following we show that $k_\tC + k_\tT \geq 4$ is necessary to cover arbitrary move sequences from an initial forest.
To this end, we look for longest move sequences from $\varepsilon^{n+1}$ and check whether the amount of rotations grows towards $4n$.

First, we notice that the move sequence $s_n = \tC(\tC\tR)^{n+1}\tT\tR^n$ with a $\tR/\tC$ ratio of $\frac{2n+1}{n+2}$ run on an initial forest produces a list-shaped tree $(ε ∙ (ε ∙ (ε ∙ \cdots)))$ with $n+1$ leaves.
The sequence $l_n = \tC\tR(\tC\tR\tT\tR\tR)^n\tC\tR\tT\tR\tT\tR^n$ started at this list reproduces it and thus can be repeated, in each iteration producing a $\tR/\tC$ ratio of $\frac{4n+3}{n+2}$.
Finally the sequence $e_n = \tC\tR(\tT\tR)^n\tT\tT$ unravels the list into a single leaf, contributing a $\tR/\tC$ ratio of $\frac{n+1}{1}$.
The sequence $a_n^m = s_n (l_n)^m e_n$ gets an accumulated ratio of
\[
  \frac{(2n+1)+(4n+3)m+(n+1)}{(n+2)+(n+2)m+1}
  = \frac{4nm+3m+3n+2}{nm + 2m + n + 3}
  = 4 - \frac{5m+n+10}{nm + 2m + n + 3}
  = 4 - \frac p q
\]
where $p = 5m + n + 10$ and $q = nm + 2m + n + 3$.
Intuitively, this approaches 4 as $m$ and $n$ go towards infinity, since $p/q > 0$ becomes arbitrary small.
Formally, we prove that for each $N > 0$ we can pick $n$ and $m$ such that $p/q \leq 1/N$.
Indeed $n = m = 6N$ fulfills $pN \leq q$ since $pN = (5 \cdot 6N + 6N + 10)N = 36N^2 + 10N < 36N^2 + 12N + 6N + 3 = q$.

% Formally, we prove that for each $N > 0$ there is $p/q < 1/N$ such that the ratio $\tC/\tR$ of for some sequence is at least $4 - p/q$.
% The sequence in question is $a_n^m$ for $n = m = 14N$ and $p = 6 \cdot 14N + 10$ and $q = (14N)^2 + 3\cdot(14N) + 3$
% Note that the $\tC/\tR$ ratio of $a_n^m = a_{14N}^{14N}$ is
% \[
%   \frac{4\cdot (14N)^2 + 6\cdot 14N + 2}{(14N)^2 + 3\cdot(14N) + 3}
%   = \frac{4q - 6\cdot(14N) - 10}{q}
%   = \frac{4q - p}{q}
%   = 4 - \frac p q
%   .
% \]
% We also have $Np = 6 \cdot 14N^2 + 10N < 14 \cdot 14N^2 + 42N + 3 = q$, thus $p/q \leq 1/N$.

%\Rotation{}



\section{Conclusion}\label{sec:conclusion}


%% Acknowledgments
% \begin{acks}                            %% acks environment is optional
                                        %% contents suppressed with 'anonymous'
  %% Commands \grantsponsor{<sponsorID>}{<name>}{<url>} and
  %% \grantnum[<url>]{<sponsorID>}{<number>} should be used to
  %% acknowledge financial support and will be used by metadata
  %% extraction tools.

% \end{acks}

%%\SHORTVERSION{\clearpage}

%% Bibliography
\bibliography{auto-main}


% %% Appendix
% \appendix
% \section{Appendix}

% Text of appendix \ldots

% \appendix


\end{document}

%%% Local Variables:
%%% ispell-local-dictionary: "british"
%%% End:
